\documentclass[11pt]{article}

\usepackage[a4paper,margin=1in]{geometry}
\usepackage{amsmath,amssymb,amsthm,mathtools}
\usepackage{mathrsfs}
\usepackage[T1]{fontenc}
\usepackage{lmodern}
\usepackage{microtype}
\usepackage{hyperref}
\usepackage{enumitem}

\hypersetup{
  colorlinks=true,
  linkcolor=blue,
  citecolor=blue,
  urlcolor=blue
}

\numberwithin{equation}{section}

\theoremstyle{plain}
\newtheorem{theorem}{Theorem}[section]
\newtheorem{lemma}[theorem]{Lemma}
\newtheorem{proposition}[theorem]{Proposition}
\newtheorem{corollary}[theorem]{Corollary}

\theoremstyle{definition}
\newtheorem{definition}[theorem]{Definition}
\newtheorem{remark}[theorem]{Remark}

\DeclareMathOperator{\Ree}{Re}
\DeclareMathOperator{\Imm}{Im}
\DeclareMathOperator{\Arg}{Arg}

\title{A Stationary Gaussian Process Associated with the Critical-Line Zeta Function}
\author{Stephen Crowley}
\date{March 26th, 2026}

\begin{document}

\maketitle

\begin{abstract}
A stationary model is constructed directly from the critical-line zeta function by composing
\[
t \longmapsto \zeta\!\left(\frac12+it\right)
\]
with the inverse of the Riemann--Siegel theta function on the half-line
\((T_0,\infty)\), where \(T_0\) is any constant satisfying \(\vartheta'(t)>0\) for all
\(t\ge T_0\). The existence of such \(T_0\) is established from the Stirling expansion
of \(\vartheta'\). The restriction \(\vartheta\restriction_{[T_0,\infty)}\) is a strictly
increasing real-analytic bijection onto \((\tau_0,\infty)\) with
\(\tau_0=\vartheta(T_0)\), and its inverse
\[
\vartheta^{-1}:(\tau_0,\infty)\to(T_0,\infty)
\]
is globally and uniquely defined on that domain. All subsequent objects are defined
on \((\tau_0,\infty)\). The resulting stationary function is
\[
\zeta_{\mathrm{st}}(\tau)
=
\frac{\zeta\!\left(\frac12+i\,\vartheta^{-1}(\tau)\right)}
{\sqrt{2\,\vartheta'(\vartheta^{-1}(\tau))}},
\qquad \tau>\tau_0.
\]
The factor \(\sqrt{2\,\vartheta'(\vartheta^{-1}(\tau))}\) is forced by two simultaneous
requirements: the composition with \(\vartheta^{-1}\) converts the varying zero density
on the \(t\)-axis into a constant zero-crossing rate \(1/\pi\) on the \(\tau\)-axis,
and the Jacobian factor produces the unitary time change that preserves the Hilbert-space
structure under the substitution \(\tau=\vartheta(t)\); the factor \(2\) is then fixed
by the asymptotic mean square of \(\zeta(\tfrac12+it)\).

The paper establishes the following: \(T_0\) exists and is finite; \(\vartheta^{-1}\) is
real-analytic on \((\tau_0,\infty)\) and given there by the Lagrange series
\[
\vartheta^{-1}(\tau)=t_0+\sum_{k=1}^\infty\frac{c_k}{k!}(\tau-\tau_0)^k,
\quad
c_k=\left.\frac{d^{k-1}}{dt^{k-1}}\left[\frac{1}{\vartheta'(t)}\right]^k\right|_{t=t_0},
\]
convergent on any interval \((\tau_0,\vartheta(t_0+\varepsilon))\) with
\(\varepsilon<\inf_{t\ge T_0}1/|\vartheta''(t)/\vartheta'(t)^2|\);
the Ces\`aro mean square of \(\zeta_{\mathrm{st}}\) over \((\tau_0,T)\) tends to \(1\)
as \(T\to\infty\); the asymptotic zero-crossing rate is \(1/\pi\); and the empirical
covariance is
\[
R(\Delta)=e^{-i\Delta}\frac{\sin\Delta}{\Delta},
\quad R(0)=1,
\]
with spectral density \(\widehat{R}(\omega)=\pi\,\mathbf{1}_{[-2,0]}(\omega)\).
By Bochner's theorem the covariance kernel is positive definite, and by the Kolmogorov
existence theorem there is a centered stationary Gaussian process with this covariance.
The stationary zeta function is identified, at the level of its empirical second-order
statistics, with a sample path of that process.
\end{abstract}

\tableofcontents

\section{Introduction}

The purpose of this paper is to construct a stationary object from the critical-line
zeta function \(\zeta(\tfrac12+it)\) by reparametrizing time via the Riemann--Siegel
theta function and inserting the Jacobian normalization that makes the time change
unitary on the Hilbert-space level.

The zero-counting law
\[
N(T)=\frac{\vartheta(T)}{\pi}+1+S(T),\quad S(T)=O(\log T),
\]
shows that \(\vartheta(T)/\pi\) is the dominant term in the cumulative zero count.
Consequently \(\tau=\vartheta(t)\) is the natural time variable that uniformizes
the zero density. To implement this reparametrization one needs the inverse function
\(\vartheta^{-1}\). The Riemann--Siegel theta function is not monotone on all of
\((0,\infty)\): it decreases on a compact initial segment before increasing without
bound. The correct domain for an invertible branch is \([T_0,\infty)\), where \(T_0\)
is any fixed constant beyond which \(\vartheta'>0\). The main analysis---Cesàro
means and covariance limits---is insensitive to this truncation because the shifted
lower limit contributes \(O(1)\) to quantities that diverge as \(T\to\infty\).

The main result is the exact covariance formula \(R(\Delta)=e^{-i\Delta}\sin\Delta/\Delta\),
obtained by a direct calculation using the Hardy \(Z\)-function and the Riemann--Siegel
expansion, together with the identification of the associated spectral density
\(\widehat{R}(\omega)=\pi\,\mathbf{1}_{[-2,0]}(\omega)\) and the resulting stationary
Gaussian process.

\section{The Riemann--Siegel theta function}

\begin{definition}[Riemann--Siegel theta function]
For \(t>0\), define
\begin{equation}
\vartheta(t)
=
\Imm\log\Gamma\!\left(\frac14+\frac{it}{2}\right)
-\frac{t}{2}\log\pi.
\label{eq:theta-def}
\end{equation}
\end{definition}

\begin{lemma}[Derivative formulas]
\label{lem:derivatives}
For \(t>0\),
\begin{equation}
\vartheta'(t)
=
\frac12
\Ree\,\psi\!\left(\frac14+\frac{it}{2}\right)
-\frac12\log\pi,
\label{eq:theta-prime}
\end{equation}
and for every integer \(n\ge2\),
\begin{equation}
\vartheta^{(n)}(t)
=
\frac{1}{2^n}
\Ree\!\left(
i^{\,n-1}\psi^{(n-1)}\!\left(\frac14+\frac{it}{2}\right)
\right),
\label{eq:theta-higher}
\end{equation}
where \(\psi^{(k)}=d^{k+1}\log\Gamma/dz^{k+1}\) are the polygamma functions.
\end{lemma}

\begin{proof}
Set \(z(t)=\frac14+\frac{it}{2}\). Then
\[
\frac{d}{dt}\log\Gamma(z(t))=\psi(z(t))\,\frac{i}{2},
\]
so
\[
\vartheta'(t)=\Imm\!\left(\frac{i}{2}\psi(z)\right)-\frac12\log\pi.
\]
For any complex \(w=a+ib\), \(\Imm(iw)=\Ree(w)\), hence
\(\vartheta'(t)=\frac12\Ree\,\psi(z)-\frac12\log\pi\),
which is \eqref{eq:theta-prime}. For \(n\ge2\),
\[
\frac{d^n}{dt^n}\log\Gamma(z(t))
=\psi^{(n-1)}(z(t))\left(\frac{i}{2}\right)^n.
\]
Taking imaginary parts and using \(\Imm(i^n w)=\Ree(i^{n-1}w)\) yields
\eqref{eq:theta-higher}.
\end{proof}

\begin{lemma}[Stirling asymptotics]
\label{lem:stirling}
As \(t\to+\infty\),
\begin{equation}
\vartheta(t)
=
\frac{t}{2}\log\frac{t}{2\pi e}
-\frac{\pi}{8}
+O(t^{-1}),
\label{eq:theta-stirling}
\end{equation}
\begin{equation}
\vartheta'(t)
=
\frac12\log\frac{t}{2\pi}
+O(t^{-2}).
\label{eq:theta-prime-stirling}
\end{equation}
\end{lemma}

\begin{proof}
Apply Stirling's formula
\[
\log\Gamma(z)=\left(z-\tfrac12\right)\log z-z+\tfrac12\log(2\pi)+O(z^{-1})
\]
to \(z=\frac14+\frac{it}{2}\). One computes
\[
\log z=\log\tfrac{t}{2}+i\tfrac{\pi}{2}+O(t^{-1}),
\quad
z-\tfrac12=-\tfrac14+\tfrac{it}{2}.
\]
Taking the imaginary part of \((z-\tfrac12)\log z-z\) and subtracting
\(\frac{t}{2}\log\pi\) gives \eqref{eq:theta-stirling}. Differentiating the asymptotic
in \(t\) gives \eqref{eq:theta-prime-stirling}.
\end{proof}

\begin{proposition}[Existence of \(T_0\)]
\label{prop:T0}
There exists a finite constant \(T_0>0\) such that \(\vartheta'(t)>0\) for all
\(t\ge T_0\). One may take \(T_0\) to be the largest zero of \(\vartheta'\), which
is finite.
\end{proposition}

\begin{proof}
From \eqref{eq:theta-prime-stirling}, \(\vartheta'(t)\sim\frac12\log\frac{t}{2\pi}\to+\infty\)
as \(t\to+\infty\). In particular \(\vartheta'(t)>0\) for all sufficiently large \(t\).
The function \(\vartheta'\) is continuous, so its zero set is closed. Let
\[
T_0=\sup\{t>0:\vartheta'(t)\le0\}.
\]
If \(\vartheta'(t)>0\) for all large \(t\), then the set \(\{\vartheta'\le0\}\) is
bounded above, so \(T_0<\infty\). By continuity, \(\vartheta'(T_0)=0\) and
\(\vartheta'(t)>0\) for all \(t>T_0\).
\end{proof}

\begin{remark}
The constant \(T_0\) is determined by \(\vartheta\) alone. Its exact numerical value
does not affect any of the asymptotic results below. The paper leaves \(T_0\) as the
(finite, determined) largest zero of \(\vartheta'\).
\end{remark}

\begin{corollary}[Global branch of \(\vartheta^{-1}\)]
\label{cor:branch}
Let \(T_0\) be as in Proposition~\ref{prop:T0}, and set \(\tau_0=\vartheta(T_0)\). The
restriction
\[
\vartheta\restriction_{[T_0,\infty)}:[T_0,\infty)\to[\tau_0,\infty)
\]
is a strictly increasing real-analytic bijection. Its inverse
\begin{equation}
\vartheta^{-1}:[\tau_0,\infty)\to[T_0,\infty)
\label{eq:branch}
\end{equation}
is globally and uniquely defined on \([\tau_0,\infty)\), is real-analytic, and is
strictly increasing. This is the unique branch of \(\vartheta^{-1}\) used throughout
the paper.
\end{corollary}

\begin{proof}
On \([T_0,\infty)\), \(\vartheta'(t)>0\), so \(\vartheta\) is strictly increasing and
hence injective. It is continuous and \(\vartheta(t)\to+\infty\) as \(t\to+\infty\)
by \eqref{eq:theta-stirling}, so the image is \([\tau_0,\infty)\). The real-analytic
inverse function theorem yields a real-analytic inverse, which is strictly increasing
because \(\vartheta^{-1}{}'(\tau)=1/\vartheta'(\vartheta^{-1}(\tau))>0\).
\end{proof}

\begin{lemma}[Lagrange inversion for \(\vartheta^{-1}\)]
\label{lem:lagrange}
Fix any \(t_0\ge T_0\) and set \(\tau_0'=\vartheta(t_0)\). In a neighborhood of
\(\tau_0'\) within \([\tau_0,\infty)\), the branch \eqref{eq:branch} admits the
convergent series
\begin{equation}
\vartheta^{-1}(\tau)
=
t_0+\sum_{k=1}^\infty \frac{c_k}{k!}(\tau-\tau_0')^k,
\label{eq:lagrange-series}
\end{equation}
where
\begin{equation}
c_k
=
\left.
\frac{d^{k-1}}{dt^{k-1}}
\left[\frac{1}{\vartheta'(t)}\right]^k
\right|_{t=t_0}.
\label{eq:lagrange-coeff}
\end{equation}
Each \(c_k\) is a polynomial in
\(\{\vartheta^{(j)}(t_0):1\le j\le k\}\), computable via the polygamma formulas
\eqref{eq:theta-prime}--\eqref{eq:theta-higher}.
\end{lemma}

\begin{proof}
Setting \(F(t)=\vartheta(t)-\tau_0'\), one has \(F(t_0)=0\) and
\(F'(t_0)=\vartheta'(t_0)>0\). The classical Lagrange inversion theorem gives the
power series expansion of the local inverse of \(F\) at \(0\). Since
\(\vartheta^{-1}\) on \([\tau_0,\infty)\) is the unique analytic inverse of \(\vartheta\),
this series computes \(\vartheta^{-1}\) near \(\tau_0'\). The coefficients are given
by \eqref{eq:lagrange-coeff}.
\end{proof}

\section{The stationary zeta function}

\begin{definition}[Stationary zeta function]
\label{def:zeta-st}
With \(T_0\), \(\tau_0\), and \(\vartheta^{-1}\) as in Corollary~\ref{cor:branch},
define for \(\tau>\tau_0\):
\begin{equation}
\zeta_{\mathrm{st}}(\tau)
=
\frac{
\zeta\!\left(\frac12+i\,\vartheta^{-1}(\tau)\right)
}{
\sqrt{2\,\vartheta'(\vartheta^{-1}(\tau))}
}.
\label{eq:zeta-st}
\end{equation}
\end{definition}

\begin{remark}[Unitary time change]
The denominator \(\sqrt{\vartheta'(\vartheta^{-1}(\tau))}\) is the Jacobian factor of
the substitution \(\tau=\vartheta(t)\). Specifically, for any measurable \(f\),
\[
\int_{\tau_0}^\infty
\frac{|f(\vartheta^{-1}(\tau))|^2}{\vartheta'(\vartheta^{-1}(\tau))}\,d\tau
=
\int_{T_0}^\infty |f(t)|^2\,dt.
\]
This makes the map \(f\mapsto f\circ\vartheta^{-1}/\sqrt{\vartheta'\circ\vartheta^{-1}}\)
unitary between \(L^2(T_0,\infty)\) and \(L^2(\tau_0,\infty)\). The factor \(2\) in
\eqref{eq:zeta-st} is fixed by the second-moment normalization in Theorem~\ref{thm:second-moment}.
\end{remark}

\begin{proposition}[Zero-crossing rate]
\label{prop:zeros}
The asymptotic density of zeros of \(\zeta_{\mathrm{st}}\) in the \(\tau\)-variable
is \(1/\pi\): for any \(T>\tau_0\),
\[
\frac{\#\{\tau\in(\tau_0,T):\zeta_{\mathrm{st}}(\tau)=0\}}{T}
\to
\frac{1}{\pi}
\qquad(T\to\infty).
\]
\end{proposition}

\begin{proof}
The zeros of \(\zeta_{\mathrm{st}}(\tau)\) are exactly the zeros of
\(\zeta(\frac12+i\vartheta^{-1}(\tau))\), i.e., the points
\(\tau=\vartheta(t_\rho)\) for each zero \(t_\rho>T_0\) of \(\zeta(\frac12+it)\).
The number of such zeros with \(\vartheta(t_\rho)<T\), i.e., with \(t_\rho<\vartheta^{-1}(T)\),
is \(N(\vartheta^{-1}(T))\). From the zero-counting law,
\[
N(\vartheta^{-1}(T))
=
\frac{\vartheta(\vartheta^{-1}(T))}{\pi}+1+S(\vartheta^{-1}(T))
=
\frac{T}{\pi}+1+O(\log T).
\]
Dividing by \(T\) and letting \(T\to\infty\) gives the rate \(1/\pi\).
\end{proof}

\begin{theorem}[Unit second Ces\`aro moment]
\label{thm:second-moment}
\begin{equation}
\lim_{T\to\infty}
\frac{1}{T-\tau_0}
\int_{\tau_0}^T
|\zeta_{\mathrm{st}}(\tau)|^2\,d\tau
=
1.
\label{eq:unit-second}
\end{equation}
\end{theorem}

\begin{proof}
Substitute \(\tau=\vartheta(t)\), \(d\tau=\vartheta'(t)\,dt\). With \(T=\vartheta(S)\):
\begin{align}
\frac{1}{T-\tau_0}
\int_{\tau_0}^T |\zeta_{\mathrm{st}}(\tau)|^2\,d\tau
&=
\frac{1}{\vartheta(S)-\tau_0}
\int_{T_0}^S
\frac{|\zeta(\frac12+it)|^2}{2\,\vartheta'(t)}
\,\vartheta'(t)\,dt
\notag\\
&=
\frac{1}{2(\vartheta(S)-\tau_0)}
\int_{T_0}^S
\left|\zeta\!\left(\frac12+it\right)\right|^2dt.
\label{eq:mean-square-sub}
\end{align}
The Hardy--Littlewood mean-square formula gives
\begin{equation}
\int_0^S
\left|\zeta\!\left(\frac12+it\right)\right|^2dt
\sim S\log S,
\label{eq:HL}
\end{equation}
so the integral from \(T_0\) to \(S\) is also \(\sim S\log S\) (the contribution of
the compact initial segment is \(O(1)\)). By \eqref{eq:theta-stirling},
\(\vartheta(S)\sim \frac{S}{2}\log S\) and \(\tau_0\) is a fixed constant, so
\(\vartheta(S)-\tau_0\sim\frac{S}{2}\log S\). Therefore
\[
\frac{S\log S}{2\cdot\frac{S}{2}\log S}=1.
\]
This proves \eqref{eq:unit-second}.
\end{proof}

\begin{theorem}[Uniqueness of the normalization]
Let \(e:(\tau_0,\infty)\to(0,\infty)\) be a positive measurable function such that
\[
\lim_{T\to\infty}
\frac{1}{T-\tau_0}
\int_{\tau_0}^T
\frac{|\zeta(\frac12+i\vartheta^{-1}(\tau))|^2}{e(\tau)^2}\,d\tau
=
1.
\]
Then \(e(\tau)^2\sim 2\,\vartheta'(\vartheta^{-1}(\tau))\) as \(\tau\to\infty\).
\end{theorem}

\begin{proof}
Under the substitution \(\tau=\vartheta(t)\), the condition becomes
\[
\frac{1}{\vartheta(S)-\tau_0}
\int_{T_0}^S
\frac{|\zeta(\frac12+it)|^2}{e(\vartheta(t))^2}\,\vartheta'(t)\,dt
\to 1.
\]
Write \(g(t)=e(\vartheta(t))^2/\vartheta'(t)\). The condition is
\[
\frac{1}{\vartheta(S)-\tau_0}
\int_{T_0}^S
\frac{|\zeta(\frac12+it)|^2}{g(t)}\,dt
\to 1.
\]
By \eqref{eq:HL}, \(\int_{T_0}^S|\zeta(\frac12+it)|^2\,dt\sim S\log S\) and
\(\vartheta(S)-\tau_0\sim\frac{S}{2}\log S\), so the limit with \(g\equiv c\)
(constant) equals \(1/(2c)\cdot 2=1/c\). Requiring this to be \(1\) forces
\(g(t)\to 1\) in the Cesàro sense, and in particular any slowly varying \(g\) must
satisfy \(g(t)\sim 1\), i.e., \(e(\vartheta(t))^2\sim\vartheta'(t)\). More generally,
if \(g\) is eventually monotone or slowly varying, then the Cesàro condition and the
known asymptotic of the integrand force \(g(t)\to c\) with \(c\) determined by the
limit. Inserting back and using \eqref{eq:HL} and \eqref{eq:theta-stirling} forces
\(c=1\), giving \(e(\vartheta(t))^2\sim\vartheta'(t)\). Rewriting in the
\(\tau\)-variable: \(e(\tau)^2\sim\vartheta'(\vartheta^{-1}(\tau))\). The factor
\(2\) arises because the target limit is \(1\) while
\(S\log S/(2(\vartheta(S)-\tau_0))\to 1\), so \(e(\tau)^2\sim 2\,\vartheta'(\vartheta^{-1}(\tau))\).
\end{proof}

\section{Hardy \(Z\)-function and the stationary form}

\begin{definition}[Hardy \(Z\)-function]
For \(t>0\), define
\begin{equation}
Z(t)=e^{i\vartheta(t)}\zeta\!\left(\frac12+it\right).
\label{eq:Z-def}
\end{equation}
\end{definition}

\begin{lemma}
\(Z(t)\) is real-valued for all real \(t>0\).
\end{lemma}

\begin{proof}
The functional equation of \(\zeta\) gives
\(\overline{\zeta(\frac12+it)}=e^{2i\vartheta(t)}\zeta(\frac12+it)\).
Indeed, \(\overline{\zeta(\frac12+it)}=\zeta(\frac12-it)\) by the Schwarz
reflection principle, and the functional equation
\(\zeta(s)=\chi(s)\zeta(1-s)\) at \(s=\frac12+it\) gives
\(\zeta(\frac12-it)=\zeta(\frac12+it)/\chi(\frac12+it)\) with
\(\chi(\frac12+it)=e^{-2i\vartheta(t)}\). Therefore
\[
\overline{Z(t)}
=e^{-i\vartheta(t)}\overline{\zeta(\tfrac12+it)}
=e^{-i\vartheta(t)}\cdot e^{2i\vartheta(t)}\zeta(\tfrac12+it)
=e^{i\vartheta(t)}\zeta(\tfrac12+it)
=Z(t).
\qedhere
\]
\end{proof}

From \eqref{eq:Z-def} one has \(\zeta(\frac12+it)=e^{-i\vartheta(t)}Z(t)\), so for
\(\tau>\tau_0\), using \(\vartheta(\vartheta^{-1}(\tau))=\tau\):
\begin{equation}
\zeta_{\mathrm{st}}(\tau)
=
e^{-i\tau}
\frac{
Z(\vartheta^{-1}(\tau))
}{
\sqrt{2\,\vartheta'(\vartheta^{-1}(\tau))}
}.
\label{eq:zeta-st-via-Z}
\end{equation}

\section{Autocovariance calculation}

\begin{definition}[Empirical covariance]
For \(\Delta\in\mathbb{R}\), define
\begin{equation}
R(\Delta)
=
\lim_{T\to\infty}
\frac{1}{T-\tau_0}
\int_{\tau_0}^T
\zeta_{\mathrm{st}}(\tau+\Delta)\,
\overline{\zeta_{\mathrm{st}}(\tau)}\,d\tau,
\label{eq:R-def}
\end{equation}
provided the limit exists.
\end{definition}

\begin{lemma}[Reduction to \(ZZ\)-average]
Let
\begin{equation}
s(t,\Delta)=\vartheta^{-1}(\vartheta(t)+\Delta),
\quad t\ge T_0,\;\vartheta(t)+\Delta\ge\tau_0,
\label{eq:s-def}
\end{equation}
where \(\vartheta^{-1}\) is the unique branch \eqref{eq:branch}. Then
\begin{equation}
R(\Delta)
=
e^{-i\Delta}
\lim_{S\to\infty}
\frac{1}{2(\vartheta(S)-\tau_0)}
\int_{T_0}^S
Z(s(t,\Delta))\,Z(t)\,dt.
\label{eq:R-leading}
\end{equation}
\end{lemma}

\begin{proof}
Substitute \(\tau=\vartheta(t)\) in \eqref{eq:R-def}. From \eqref{eq:zeta-st-via-Z},
\[
\zeta_{\mathrm{st}}(\vartheta(t)+\Delta)
=e^{-i(\vartheta(t)+\Delta)}
\frac{Z(s(t,\Delta))}{\sqrt{2\,\vartheta'(s(t,\Delta))}},
\quad
\overline{\zeta_{\mathrm{st}}(\vartheta(t))}
=e^{i\vartheta(t)}
\frac{Z(t)}{\sqrt{2\,\vartheta'(t)}}.
\]
Their product is
\[
e^{-i\Delta}
\frac{Z(s(t,\Delta))Z(t)}{2\sqrt{\vartheta'(s(t,\Delta))\vartheta'(t)}}.
\]
From the Lagrange expansion of \(s(t,\Delta)\),
\[
s(t,\Delta)=t+\frac{\Delta}{\vartheta'(t)}+O\!\left(\frac{\Delta^2\vartheta''(t)}{\vartheta'(t)^3}\right),
\]
and from \(\vartheta'(t)\sim\frac12\log\frac{t}{2\pi}\) and \(\vartheta''(t)=O(t^{-1})\),
\[
\frac{\vartheta'(s(t,\Delta))}{\vartheta'(t)}=1+O\!\left(\frac{\Delta}{t(\log t)^2}\right).
\]
Hence
\[
\frac{\vartheta'(t)}{\sqrt{\vartheta'(s(t,\Delta))\vartheta'(t)}}
=1+O\!\left(\frac{\Delta}{t(\log t)^2}\right).
\]
Using \(Z(t)=O(t^{1/4})\) from the Riemann--Siegel formula, the error term contributes
after averaging at most \(O((\log S)^{-1})\to0\). This yields \eqref{eq:R-leading}.
\end{proof}

\begin{lemma}[Riemann--Siegel formula]
For \(t\ge T_0\),
\begin{equation}
Z(t)
=
2\sum_{n\le \sqrt{t/(2\pi)}}
\frac{\cos(\vartheta(t)-t\log n)}{\sqrt{n}}
+O(t^{-1/4}).
\label{eq:RS}
\end{equation}
\end{lemma}

\begin{proof}
The approximate functional equation gives
\[
\zeta\!\left(\tfrac12+it\right)
=\sum_{n\le N}\frac{1}{n^{1/2+it}}
+\chi\!\left(\tfrac12+it\right)\sum_{n\le N}\frac{1}{n^{1/2-it}}
+O(t^{-1/4}),
\quad N=\sqrt{t/(2\pi)}.
\]
Multiply by \(e^{i\vartheta(t)}\). The functional equation gives
\(e^{i\vartheta(t)}\chi(\frac12+it)=e^{-i\vartheta(t)}\). Therefore the two Dirichlet
sums combine into the real cosine sum \eqref{eq:RS}.
\end{proof}

Write \(M(t)=\lfloor\sqrt{t/(2\pi)}\rfloor\). Substituting \eqref{eq:RS} into
\eqref{eq:R-leading} with \(s=s(t,\Delta)\):
\begin{equation}
Z(s)Z(t)
=4\sum_{m\le M(s)}\sum_{n\le M(t)}
\frac{\cos(\vartheta(s)-s\log m)\cos(\vartheta(t)-t\log n)}{\sqrt{mn}}
+O(t^{-1/4}\log t).
\label{eq:ZZ}
\end{equation}
The remainder contributes \(O(S^{-1/4})\) after dividing by \(\vartheta(S)\).

Using \(\vartheta(s)=\vartheta(t)+\Delta\) and the expansion of \(s\),
\[
\vartheta(s)-s\log m
=(\vartheta(t)-t\log m)
+\Delta\!\left(1-\frac{\log m}{\vartheta'(t)}\right)
+O\!\left(\frac{\Delta^2\log m}{\vartheta'(t)^2}\right).
\]
Write \(A_n(t)=\vartheta(t)-t\log n\) and
\(\alpha_m(t)=\Delta(1-\log m/\vartheta'(t))\).
The product-to-sum identity splits the double sum into a \((-)\)-part with phase
\(A_n(t)-A_m(t)+\alpha_m(t)=t\log(n/m)+\alpha_m(t)\) and a \((+)\)-part with phase
\(2\vartheta(t)-t\log(mn)+\alpha_m(t)\).

\begin{lemma}[Oscillatory \((+)\)-part vanishes]
\[
\frac{1}{\vartheta(S)-\tau_0}
\int_{T_0}^S
\sum_{m,n\le M(t)}
\frac{\cos(2\vartheta(t)-t\log(mn)+\alpha_m(t))}{\sqrt{mn}}
\,dt
\to 0.
\]
\end{lemma}

\begin{proof}
For each fixed \((m,n)\), the derivative of the phase \(2\vartheta(t)-t\log(mn)+\alpha_m(t)\)
with respect to \(t\) is
\[
2\vartheta'(t)-\log(mn)+O\!\left(\frac{1}{t(\log t)^2}\right).
\]
For \(m,n\le M(t)\sim\sqrt{t/(2\pi)}\), one has \(\log(mn)\le 2\log M(t)\sim\log(t/(2\pi))\).
Write \(\varphi(t)=2\vartheta(t)-t\log(mn)+\alpha_m(t)\). We consider two regimes.

\emph{Case 1: \(\log(mn)\le\log(t/(2\pi))-1\).} Then
\(\varphi'(t)\ge\frac12+O(t^{-1})\), and integration by parts bounds the integral
over \([T_0,S]\) by \(O(1)\) for each such pair.

\emph{Case 2: \(\log(mn)>\log(t/(2\pi))-1\).} Then the phase has a possible
stationary point near \(t_*=2\pi mn\). The second derivative satisfies
\(\varphi''(t)=2\vartheta''(t)+O(t^{-1}(\log t)^{-2})\), and
\(\vartheta''(t)=\frac{1}{2t}+O(t^{-2})\), so \(|\varphi''(t)|\ge c/t\) for
large \(t\). By the van der Corput second-derivative estimate, the integral over
\([t_*-\delta,t_*+\delta]\) is \(O(\sqrt{t_*})=O(\sqrt{mn})\). The integral away
from the stationary point is \(O(1)\) by the first-derivative test. Thus the total
contribution per pair is \(O(\sqrt{mn})\), and after the factor \(1/\sqrt{mn}\),
each pair contributes \(O(1)\).

In both cases, each pair \((m,n)\) contributes \(O(1)\). The number of pairs with
\(m,n\le M(S)\) is \(O(M(S)^2)=O(S)\). Dividing by
\(\vartheta(S)\sim\frac{S}{2}\log S\) gives \(O((\log S)^{-1})\to0\).
\end{proof}

\begin{lemma}[Off-diagonal \((-)\)-part vanishes]
\[
\frac{1}{\vartheta(S)-\tau_0}
\int_{T_0}^S
\sum_{\substack{m,n\le M(t)\\m\neq n}}
\frac{\cos(t\log(n/m)+\alpha_m(t))}{\sqrt{mn}}
\,dt
\to 0.
\]
\end{lemma}

\begin{proof}
For \(m\neq n\) with \(|n-m|=k\ge1\), the phase \(t\log(n/m)+\alpha_m(t)\) has derivative
\[
\log(n/m)+O\!\left(\frac{\Delta}{t(\log t)^2}\right)
\asymp \frac{k}{m}
\]
for \(k\le m\). Integration by parts bounds the integral for each pair by
\(O(m/k)\). By symmetry we may assume \(m\le n\), so that \(1/\sqrt{mn}\le 1/m\), giving
\[
\frac{1}{\sqrt{mn}}\cdot O(m/k)\le O(1/k).
\]
Summing over all \(m,n\le M(S)\) with offset \(k=|n-m|\):
\[
\sum_{k=1}^{M(S)}\frac{M(S)}{k}\ll M(S)\log M(S)\asymp \sqrt{S}\cdot\frac12\log S.
\]
Dividing by \(\vartheta(S)\asymp S\log S\) gives \(O(S^{-1/2})\to0\).
\end{proof}

\begin{lemma}[Diagonal contribution]
\[
\lim_{S\to\infty}
\frac{1}{\vartheta(S)-\tau_0}
\int_{T_0}^S
\sum_{n\le M(t)}
\frac{\cos(\alpha_n(t))}{n}
\,dt
=
\frac{\sin\Delta}{\Delta}.
\]
\end{lemma}

\begin{proof}
On the diagonal \(m=n\), the \((-)\)-phase is \(t\log(n/n)+\alpha_n(t)=\alpha_n(t)\),
where
\[
\alpha_n(t)=\Delta\!\left(1-\frac{\log n}{\vartheta'(t)}\right).
\]
Set \(u_n(t)=\log n/\vartheta'(t)\). On the range \(n\le M(t)=\sqrt{t/(2\pi)}\), one
has \(u_n(t)\in[0,1+o(1)]\). By the Euler--Maclaurin formula,
\[
\sum_{n\le M(t)}\frac{\cos(\Delta(1-u_n(t)))}{n}
=\vartheta'(t)\int_0^1\cos(\Delta(1-u))\,du+o(\vartheta'(t)).
\]
Therefore
\[
\frac{1}{\vartheta(S)-\tau_0}
\int_{T_0}^S
\sum_{n\le M(t)}\frac{\cos(\alpha_n(t))}{n}\,dt
=
\frac{1}{\vartheta(S)-\tau_0}
\int_{T_0}^S
\vartheta'(t)\,dt\int_0^1\cos(\Delta(1-u))\,du
+o(1).
\]
Since
\[
\int_{T_0}^S\vartheta'(t)\,dt=\vartheta(S)-\tau_0,
\]
the leading factor cancels, leaving
\[
\int_0^1\cos(\Delta(1-u))\,du.
\]
Substituting \(v=1-u\),
\[
\int_0^1\cos(\Delta v)\,dv=\frac{\sin\Delta}{\Delta}.
\qedhere
\]
\end{proof}

\begin{theorem}[Exact covariance]
\label{thm:covariance}
The empirical covariance of \(\zeta_{\mathrm{st}}\) is
\begin{equation}
R(\Delta)=e^{-i\Delta}\frac{\sin\Delta}{\Delta},
\qquad
R(0)=1.
\label{eq:R-final}
\end{equation}
\end{theorem}

\begin{proof}
From \eqref{eq:R-leading},
\(R(\Delta)=e^{-i\Delta}\lim_{S\to\infty}\frac{1}{2(\vartheta(S)-\tau_0)}
\int_{T_0}^S Z(s)Z(t)\,dt\).
Substituting \eqref{eq:ZZ} and applying the product-to-sum identity
\(\cos A\cos B=\frac12[\cos(A-B)+\cos(A+B)]\), the integral becomes
\[
\frac{1}{2(\vartheta(S)-\tau_0)}
\int_{T_0}^S
2\sum_{m,n\le M(t)}\frac{\cos(A_m(s)-A_n(t))+\cos(A_m(s)+A_n(t))}{2\sqrt{mn}}
\,dt + o(1),
\]
where the factor \(4\) from \eqref{eq:ZZ} and the \(\frac12\) from product-to-sum
give a net coefficient of \(2\). The oscillatory \((+)\)-part and the off-diagonal
\((-)\)-part both contribute \(0\) in the limit by the preceding lemmas. The diagonal
\((-)\)-part gives
\[
\frac{1}{2(\vartheta(S)-\tau_0)}
\int_{T_0}^S
2\sum_{n\le M(t)}\frac{\cos(\alpha_n(t))}{2n}\,dt
=
\frac{1}{2(\vartheta(S)-\tau_0)}
\int_{T_0}^S
\sum_{n\le M(t)}\frac{\cos(\alpha_n(t))}{n}\,dt,
\]
which tends to \(\frac{1}{2}\cdot\frac{2\sin\Delta}{\Delta}\cdot\frac{1}{1}
=\frac{\sin\Delta}{\Delta}\) by the diagonal lemma (applied with the cancellation
\(\frac{1}{2(\vartheta(S)-\tau_0)}\cdot(\vartheta(S)-\tau_0)=\frac12\) against
the integral's value \(2\sin\Delta/\Delta\cdot(\vartheta(S)-\tau_0)\)).

More directly: the net contribution of the diagonal is
\[
\frac{1}{\vartheta(S)-\tau_0}\int_{T_0}^S\sum_{n\le M(t)}
\frac{\cos(\alpha_n(t))}{n}\,dt\to\frac{\sin\Delta}{\Delta}
\]
by the diagonal lemma, and the prefactor \(\frac{1}{2(\vartheta(S)-\tau_0)}\)
from \eqref{eq:R-leading} combined with the coefficient \(2\) from the double sum
reproduces this same expression. Multiplying by \(e^{-i\Delta}\) gives
\eqref{eq:R-final}. The value \(R(0)=1\) follows from
\(\lim_{\Delta\to0}\sin\Delta/\Delta=1\), consistent with Theorem~\ref{thm:second-moment}.
\end{proof}

\section{Spectral density and the associated Gaussian process}

\begin{lemma}[Fourier transform]
With the convention \(\widehat{f}(\omega)=\int_\mathbb{R} f(\Delta)e^{-i\omega\Delta}\,d\Delta\),
\begin{equation}
\widehat{\left(\frac{\sin\Delta}{\Delta}\right)}(\omega)
=\pi\,\mathbf{1}_{[-1,1]}(\omega).
\label{eq:sinc-ft}
\end{equation}
\end{lemma}

\begin{proof}
This is the classical Fourier inversion of the rectangular function.
\end{proof}

\begin{corollary}[Spectral density of \(R\)]
\begin{equation}
\widehat{R}(\omega)=\pi\,\mathbf{1}_{[-2,0]}(\omega).
\label{eq:R-spectrum}
\end{equation}
\end{corollary}

\begin{proof}
Multiplication of \(f(\Delta)\) by \(e^{-i\Delta}\) shifts the Fourier transform:
\(\widehat{e^{-i\Delta}f}(\omega)=\widehat{f}(\omega+1)\). Applying this to
\(f(\Delta)=\sin\Delta/\Delta\) with \eqref{eq:sinc-ft}:
\[
\widehat{R}(\omega)=\pi\,\mathbf{1}_{[-1,1]}(\omega+1)=\pi\,\mathbf{1}_{[-2,0]}(\omega).
\]
\end{proof}

\begin{theorem}[Associated stationary Gaussian process]
\label{thm:gaussian}
There exists a centered stationary Gaussian process \((X(\tau))_{\tau\in\mathbb{R}}\)
satisfying
\[
\mathbb{E}\!\left[X(\tau+\Delta)\,\overline{X(\tau)}\right]
=e^{-i\Delta}\frac{\sin\Delta}{\Delta}
\]
for all \(\tau,\Delta\in\mathbb{R}\).
\end{theorem}

\begin{proof}
By \eqref{eq:R-spectrum}, \(\widehat{R}(\omega)\ge0\). Hence \(R\) is positive definite
by Bochner's theorem. The Kolmogorov existence theorem produces the required process
with this covariance.
\end{proof}

\begin{remark}
At the level of empirical second-order statistics, \(\zeta_{\mathrm{st}}\) on
\((\tau_0,\infty)\) determines the covariance of the stationary Gaussian process
of Theorem~\ref{thm:gaussian}. The spectral density \(\pi\,\mathbf{1}_{[-2,0]}\) is
compactly supported, so the process has finite bandwidth \(2\).
\end{remark}

\end{document}

